Para corregir la interpretación de la salida sigmoide, se ajustó un calibrador post-hoc usando solamente el conjunto de validación. Se compararon Platt e isotónica por \emph{Brier score} de validación y se seleccionó \textbf{isotonic}. El test se mantuvo como diagnóstico congelado: no se usó para elegir el calibrador.

En test, la salida sigmoide cruda obtuvo Brier 0.2017 y ECE 0.3171; el score calibrado obtuvo Brier 0.0970 y ECE 0.0169. Por eso, la GUI deja de presentar la salida como ``probabilidad operacional'' y la muestra como \textbf{score calibrado de riesgo mortal}.

Además, se calcularon intervalos de confianza bootstrap al 95\,\% sobre las predicciones congeladas del test: F1-mortal 0.2941 [0.2441, 0.3482], recall-mortal 0.4895 [0.4107, 0.5724] y PR-AUC 0.2890 [0.2244, 0.3727]. Estos intervalos evitan vender el resultado como un punto exacto: muestran la incertidumbre real de un test pequeño y desbalanceado.
